\documentclass[aspectratio=169]{beamer}
\usetheme{Madrid}
\usecolortheme{dolphin}
\setbeamertemplate{navigation symbols}{}
\setbeamertemplate{footline}[frame number]
\setbeamersize{text margin left=7mm,text margin right=7mm}
\usepackage[T1]{fontenc}
\usepackage{lmodern}
\usepackage{booktabs,tabularx,array}
\usepackage{xcolor}
\hypersetup{hidelinks,pdftitle={Mixed-radix: Final Evidence and Closure}}
\newcommand{\good}[1]{\textcolor{green!45!black}{\textbf{#1}}}
\newcommand{\bad}[1]{\textcolor{red!75!black}{\textbf{#1}}}
\title{Original Mixed-radix: Final Evidence and Closure}
\subtitle{A real mechanism, but no stable natural-data method}
\author{}
\date{2026-08-04}

\begin{document}
\begin{frame}\titlepage\end{frame}

\begin{frame}{1. The Answer First}
\begin{alertblock}{Final result for the tested method}
The original fixed-adjacent mixed-radix method is \textbf{not competitive on
SIFT1M or GIST1M} under the frozen Recall--speed evaluation.
\end{alertblock}
\begin{itemize}
  \item At 32 bytes, it becomes exactly the power-of-two control.
  \item At 64 bytes, Recall changes by at most about one tenth of one
  percentage point, with both positive and negative signs.
  \item A synthetic positive control now confirms that the intended mechanism
  itself is real and reaches the query result.
\end{itemize}
\begin{block}{Interpretation}
The implementation can exploit non-power-of-two cardinalities. The tested
natural data simply does not contain enough of this structure for a useful win.
\end{block}
\end{frame}

\begin{frame}{2. Scope Correction: Which ``Attempt 4'' Is This?}
\begin{columns}[T]
\column{0.47\textwidth}
\begin{block}{Original target: this deck}
Two adjacent scalar quantizers share one packed word. Their numbers of levels
may be arbitrary positive integers.
\end{block}
\column{0.47\textwidth}
\begin{block}{Separate branch: not this result}
The structured-2D method learns a reusable two-dimensional shape with affine
transforms. It was evaluated separately and is currently set aside.
\end{block}
\end{columns}
\begin{alertblock}{Why the distinction matters}
Mixed-radix changes how a fixed number of bit patterns is divided between two
scalar coordinates. It does \textbf{not} learn a shared 2D vector codebook.
\end{alertblock}
\end{frame}

\begin{frame}{3. The Original Intuition}
Four bits give exactly \(2^4=16\) possible labels for one coordinate pair.
\medskip
\begin{columns}[T]
\column{0.47\textwidth}
\begin{block}{Power-of-two allocation}
Each coordinate receives 1, 2, 4, 8, \ldots levels. A typical balanced split
is \(4\times4=16\).
\end{block}
\column{0.47\textwidth}
\begin{block}{Mixed-radix allocation}
Allow any integers \(K_1,K_2\) as long as
\[
K_1K_2\leq16.
\]
This includes \(3\times5=15\).
\end{block}
\end{columns}
\begin{alertblock}{Hypothesis}
If one coordinate really needs 3 values and the other needs 5, mixed-radix can
retain all 15 combinations. Separate binary rounding would need
\(4\times8=32\) states and cannot fit in four bits.
\end{alertblock}
\end{frame}

\begin{frame}{4. Running Example: Pack a \(3\times5\) Grid}
\small
Let \(z_1\in\{0,1,2\}\) and \(z_2\in\{0,1,2,3,4\}\). Store
\[
\text{label}=z_1+3z_2.
\]
\begin{center}
\begin{tabular}{c|ccccc}
\toprule
& \(z_2=0\) & 1 & 2 & 3 & 4\\
\midrule
\(z_1=0\) & 0 & 3 & 6 & 9 & 12\\
\(z_1=1\) & 1 & 4 & 7 & 10 & 13\\
\(z_1=2\) & 2 & 5 & 8 & 11 & 14\\
\bottomrule
\end{tabular}
\end{center}
Every real combination has a unique label; label 15 is simply unused.
\medskip
\begin{block}{What the power-of-two control must do}
With only 16 labels it may choose \(4\times4\), but then five distinct values
of the second coordinate must share only four reconstructed values. At least
two real prototypes can receive the same stored code.
\end{block}
\end{frame}

\begin{frame}[fragile]{5. Training, Encoding, and Querying}
\small
\begin{verbatim}
Training, using base/index data only:
  for each adjacent coordinate pair:
    estimate scalar reconstruction cost for each level count
    choose K1, K2 minimizing total cost, subject to K1*K2 <= 2^B

Encoding one database vector:
  z1 = nearest scalar centre on coordinate 1
  z2 = nearest scalar centre on coordinate 2
  store label = z1 + K1*z2

Querying:
  build the complete 2^B-entry distance table for each pair
  decode one label per pair, add table values, keep the best 100 IDs
\end{verbatim}
\begin{alertblock}{Controlled difference}
Arbitrary arm A and power-of-two arm D use the same pairs, bytes, candidates,
table builder, and scan loop. Only the allowed values of \(K_1,K_2\) differ.
\end{alertblock}
\end{frame}

\begin{frame}{6. What We Actually Compared}
\small
\begin{tabularx}{\textwidth}{>{\bfseries}p{0.26\textwidth} X}
A128 & arbitrary integer radices, such as \(15\times17=255\)\\
D128 FULL & identical method, but each radix must be a power of two\\
PQ / OPQ & standard product-quantization baselines for deployability context\\
data & SIFT1M and GIST1M\\
payload & exactly 32 or 64 code bytes per database vector\\
index sizes & 1,024 and 4,096 coarse database cells\\
quality & Recall@100: how many of the true nearest 100 are returned\\
speed & single-query latency and 12-core queries per second\\
search range & nine fixed numbers of cells searched per query\\
\end{tabularx}
\begin{block}{Fairness rule}
A and D search exactly the same candidates and perform the same table work.
The allocation was fixed before reading query results.
\end{block}
\end{frame}

\begin{frame}{7. What the Natural-data Allocator Chose}
\begin{columns}[T]
\column{0.47\textwidth}
\begin{block}{32-byte vectors: 4 bits per pair}
Every A group selected \(4\times4\). Therefore A and D were literally the
same representation and returned identical rankings.
\end{block}
\column{0.47\textwidth}
\begin{block}{64-byte vectors: 8 bits per pair}
Only 10--24 of 64 groups differed from \(16\times16\). Almost all were
\(15\times17\) or \(17\times15\), using 255 of 256 paid states.
\end{block}
\end{columns}
\begin{alertblock}{Early mechanism diagnosis}
The reconstruction-based allocator found no large unused-cardinality problem
on SIFT/GIST. At 64 bytes it mostly moved only one level from one coordinate
to its partner.
\end{alertblock}
\end{frame}

\begin{frame}{8. Natural-data Recall Result}
\small
At the same number of searched cells, A-minus-D Recall@100 was:
\begin{center}
\begin{tabular}{lrrr}
\toprule
Dataset & coarse cells & minimum & maximum\\
\midrule
GIST1M & 1,024 & -0.00096 & +0.00054\\
GIST1M & 4,096 & -0.00106 & +0.00017\\
SIFT1M & 1,024 & -0.000097 & +0.000066\\
SIFT1M & 4,096 & -0.000033 & +0.000436\\
\bottomrule
\end{tabular}
\end{center}
\begin{alertblock}{Meaning in plain language}
The largest gain was 0.054 percentage points; the largest loss was 0.106
percentage points. The sign changed across datasets, index sizes, and search
ranges, so there is no stable quality improvement.
\end{alertblock}
\end{frame}

\begin{frame}{9. Speed and Baseline Result}
\begin{itemize}
  \item A/D throughput ratio: \(0.9918\) to \(1.0069\); group medians are
  essentially 1.
  \item A/D p95 latency ratio: \(0.9897\) to \(1.0251\).
  \item This is expected: A and D deliberately execute the same complete-table
  consumer and nearly the same representation.
  \item PQ and OPQ generally reach higher Recall in the useful high-quality
  region.
\end{itemize}
\begin{alertblock}{Natural-data decision}
The fixed-adjacent mixed-radix method did not move a useful Recall--speed
frontier. A tiny nonzero difference is not a database-systems contribution.
\end{alertblock}
\scriptsize Complete matrix: 512 accepted passes; 32.36 CPU-hours;
17.09 wall-hours; all eight correctness tests pass.
\end{frame}

\begin{frame}{10. Why Add a Synthetic Positive Control?}
The natural result left two different explanations:
\begin{enumerate}
  \item the mixed-radix mechanism is valid, but the required data structure is
  absent or weak; or
  \item the implementation cannot transmit a radix advantage into query
  ranking, even when such an advantage exists.
\end{enumerate}
\begin{block}{Decisive construction}
Create exact Cartesian prototypes with the cardinalities required by the
intuition. Put 100 identical database vectors at every prototype, place one
query exactly on each prototype, search every candidate, and ask whether all
100 true duplicates are recovered.
\end{block}
Positive cases are \(3\times5\) in 4 bits and \(15\times17\) in 8 bits.
Matching \(4\times4\) and \(16\times16\) cases are zero-effect controls.
\end{frame}

\begin{frame}{11. Synthetic Result: The Mechanism Reaches Recall}
\small
\begin{center}
\begin{tabular}{lcccc}
\toprule
Case & A shape & D shape & code collisions A/D & Recall A/D\\
\midrule
positive B4 & \(3\times5\) & \(4\times4\) & 0 / 3 & \good{1.000 / 0.800}\\
null B4 & \(4\times4\) & \(4\times4\) & 0 / 0 & 1.000 / 1.000\\
positive B8 & \(15\times17\) & \(16\times16\) & 0 / 15 & \good{1.000 / 0.941}\\
null B8 & \(16\times16\) & \(16\times16\) & 0 / 0 & 1.000 / 1.000\\
\bottomrule
\end{tabular}
\end{center}
\begin{block}{Validity checks}
In both null controls A and D produced identical rankings. In both positive
cases the rankings differed exactly where D merged prototypes. Two complete
runs were byte-for-byte identical.
\end{block}
\end{frame}

\begin{frame}{12. Put the Two Results Together}
\begin{center}
\begin{tabularx}{0.94\textwidth}{>{\bfseries}p{0.24\textwidth} X X}
\toprule
Question & Natural SIFT/GIST & Synthetic witness\\
\midrule
Can A choose non-dyadic radices? & rarely, and only weakly at 64 bytes & yes,
exactly \(3\times5\) and \(15\times17\)\\
Does this avoid collisions? & no material measured benefit & yes, 3 and 15
control collisions removed\\
Does it improve Recall? & tiny, mixed-sign changes & yes, by 0.200 and 0.0588\\
Does it beat PQ/OPQ? & no useful frontier movement & not tested or claimed\\
\bottomrule
\end{tabularx}
\end{center}
\begin{alertblock}{Combined conclusion}
The mechanism is real but its necessary structure is not sufficiently strong
in the tested natural representation. The synthetic result explains the
negative result; it does not reverse it.
\end{alertblock}
\end{frame}

\begin{frame}{13. Coordinate Regrouping Was Not the Missing Contribution}
\small
Allowing non-adjacent coordinate pairs improved Recall substantially over the
fixed pairing \((0,1),(2,3),\ldots\).  We then compared our minimum-cost
perfect matching with OPQ's existing variance-based grouping rule.
\begin{center}
\begin{tabular}{lrr}
\toprule
& search 64 cells & search 1,024 cells\\
\midrule
SIFT1M: matching minus OPQ grouping & +0.000253 & +0.000349\\
GIST1M: matching minus OPQ grouping & +0.000320 & -0.000160\\
\bottomrule
\end{tabular}
\end{center}
\begin{alertblock}{Interpretation}
``Do not pair coordinates mechanically'' is useful, but already explained by
known transform-coding practice.  Arbitrary integer radices still add at most
about 0.00057 Recall over the identical power-of-two pairing.
\end{alertblock}
\end{frame}

\begin{frame}{14. Why Reconstruction Error Was Not Enough}
\footnotesize
Recall depends on the ordering of complete estimated distances, not on the
average reconstruction error of one coordinate or segment.
\begin{block}{The diagnostic question}
For each database target, replace one stored segment contribution by its exact
distance while leaving every other estimated contribution unchanged.  How
many wrong pairwise orderings disappear?
\end{block}
\begin{itemize}
  \item This is an \textbf{oracle upper bound}: it uses information that a
  deployable compressed code does not store.
  \item A single segment repaired only about 43--44\% of inversions and missed
  the frozen absolute-improvement rule.
  \item Replacing the 192-dimensional, 6-bit segment together with the
  320-dimensional, 4-bit segment repaired \good{71.1--75.6\%}.
\end{itemize}
\begin{alertblock}{What this did and did not establish}
Two segments jointly contain ranking headroom.  Exact replacement is not an
encoder, and the nonlinear Recall ordering test can make independent errors
appear jointly important through cancellation.
\end{alertblock}
\end{frame}

\begin{frame}{15. Can a Query-unaware Joint Encoder Use That Headroom?}
\small
After the SAQ plan is fixed, the two segments' legal codes form a Cartesian
product.  Reconstruction error and unseen-direction worst-case error separate
into one objective per segment.  Coupling requires a chosen direction
distribution.
\begin{block}{Illustrative cancellation}
On training directions, two segment errors might be (+2) and (-2), so the
total error is zero.  On different directions they might become (+2) and
(+2), so the same jointly selected pair has total error four.  A cancellation
seen during training need not survive new directions.
\end{block}
\begin{alertblock}{Falsifiable hypothesis}
Choose the two codes jointly using one set of base-derived directions, freeze
the choices, and evaluate on a disjoint set.  Require at least 5\% lower MSE
than independent selection in both directions and improvement on at least
60\% of targets.
\end{alertblock}
\end{frame}

\begin{frame}[fragile]{16. Frozen Base-only Cross-fit Test}
\small
No benchmark query or ground truth was read.
\begin{verbatim}
for each of 32 held-out database targets:
  alternatives = stored code + every legal one-coordinate +/-1 move
  train on base directions 0..127:
    IND   = best code for each segment separately
    JOINT = best pair over the Cartesian product
  freeze both choices; evaluate on directions 128..255
repeat with the two direction folds swapped
\end{verbatim}
Both arms use the same alternatives, bytes, analytical rescale, and estimator.
The evaluation oracle chooses after seeing the held-out directions and is
reported only as a ceiling, never as a method.
\end{frame}

\begin{frame}{17. Joint Selection Did Not Generalize}
\small
\begin{center}
\begin{tabular}{lrr}
\toprule
& train fold 0 & train fold 1\\
\midrule
joint MSE change vs independent & +0.0805\% & \bad{-1.3644\%}\\
targets improved & 15/32 & 13/32\\
joint and independent choices differ & 31/32 & 29/32\\
held-out oracle improvement & 16.83\% & 15.92\%\\
\bottomrule
\end{tabular}
\end{center}
Positive MSE change means an improvement; the negative value means the joint
choice is worse.
\begin{alertblock}{Decision: JOINT\_LOCAL\_NO\_GO}
Good code pairs exist after seeing evaluation directions, but a pair learned
from disjoint base directions does not predict them.  The observed headroom is
direction-specific cancellation, not a stable query-unaware rule.
\end{alertblock}
\end{frame}

\begin{frame}{18. Final Claim Boundary}
\begin{block}{What remains valid}
Mixed-radix packing can preserve non-power-of-two Cartesian cardinalities and
prevent collisions in data deliberately having that structure.
\end{block}
\begin{alertblock}{What the natural-data evidence rejects}
Neither fixed adjacent mixed-radix, optimized coordinate pairing, nor the
tested base-trained two-segment joint objective produces a stable new
Recall--speed method beyond the closest controls.
\end{alertblock}
\begin{block}{Research decision}
Close this direct line.  Do not enlarge the local code neighbourhood, tune the
folds or threshold, or inspect benchmark queries to rescue it.  Continuing
would require a genuinely new question---for example changing the estimator
or storing cross-segment information---followed by a new closest-work and
cost review.
\end{block}
\end{frame}

\begin{frame}{References and Evidence}
\tiny
\begin{itemize}
  \item J\'egou, Douze, and Schmid. ``Product Quantization for Nearest
  Neighbor Search.'' TPAMI 2011.
  \item Ge et al. ``Optimized Product Quantization.'' CVPR 2013.
  \item Gao and Long. ``RaBitQ.'' SIGMOD 2024.
  \item Li et al. ``SAQ.'' 2025.
  \item Frozen mixed-radix design:
  \texttt{mixed\_radix\_query\_design\_2026\_07\_30.md}.
  \item Natural-query result:
  \texttt{mixed\_radix\_query\_results\_2026\_08\_01.md}.
  \item Synthetic mechanism result:
  \texttt{mixed\_radix\_synthetic\_witness\_result\_2026\_08\_01.md}.
  \item Closest grouping baseline:
  \texttt{nonadjacent\_pairing\_closest\_baseline\_result\_2026\_08\_03.md}.
  \item Exact component and joint-component oracle results:
  \texttt{saq\_component\_oracle\_diagnostic\_result\_2026\_08\_03.md};
  \texttt{saq\_joint\_component\_oracle\_result\_2026\_08\_03.md}.
  \item Static objective audit and base-trained cross-fit result:
  \texttt{saq\_two\_segment\_feasible\_objective\_static\_audit\_2026\_08\_03.md};
  \texttt{saq\_joint\_objective\_viability\_result\_2026\_08\_04.md}.
\end{itemize}
\end{frame}

\end{document}
